\documentclass[12pt,letterpaper]{report}
\usepackage{graphicx}
\usepackage{scrextend}
\usepackage{vmargin}
\usepackage{graphicx}
\usepackage{multirow}
\usepackage[utf8]{inputenc}
\usepackage[spanish]{babel}
\usepackage{multicol}
\usepackage{enumerate}
\usepackage{float}
\usepackage{amsmath, amsthm, amssymb, amsfonts}
\usepackage[usenames]{color}
\parindent=0mm
\pagestyle{empty}
\definecolor{miorange}{rgb}{0.91, 0.43, 0.0}
\begin{document}
\setmargins{2.5cm}      
{1.5cm}                     
{2cm}  
{24cm}                    
{10pt}                          
{1cm}                          
{0pt}                             
{2cm}
\begin{flushright}
\today
\end{flushright}
\vspace{-1.75cm}
\section*{Tarea Clase 1}
\subsection*{Efectúa las siguientes operaciones:}
\begin{enumerate}
    \item \begin{equation*}
    (5)-(7)-(-12)+(-8)
    \end{equation*}
        \item \begin{equation*}
        -12-(-3)+(-8)+(14)-(20)
    \end{equation*}
        \item \begin{equation*}
        5-[7-2-(1-9)-3+12]+4
    \end{equation*}
    \item \begin{equation*}
     6-(-9+7-1)-[3-(-5+4+6)-1]
    \end{equation*}
    \item \begin{equation*}
        28-[-21-(12-3)-7]
    \end{equation*}
    \item \begin{equation*}
        [-2+3 \cdot (2-5) \div 3] [(3-5+2)-2 \cdot (3-4)]
    \end{equation*}
    \item \begin{equation*}
        4 \cdot 2^3-3 \cdot 3 +5 \cdot (4+3)
    \end{equation*}
    \item \begin{equation*}
        4 \cdot [3+2 \cdot (4-2)+4]
    \end{equation*}
    \item \begin{equation*}
        6(-1)^3-4
    \end{equation*}
    \item \begin{equation*}
        -4-(-3)^2+\sqrt{9}
    \end{equation*}
    \item \begin{equation*}
    \frac{5}{12}- \frac{7}{9} \div (\frac{4}{3}+2)
    \end{equation*}
    \item \begin{equation*}
        \left[\frac{6}{5}\div \frac{9}{10}-\left(  2-\frac{7}{12}\right) \right] +\frac{7}{24}
    \end{equation*}
    \item \begin{equation*}
     \frac{5}{36}-\left(  \frac{7}{16}+\frac{1}{4}\div \frac{3}{5}\right)
    \end{equation*}
\end{enumerate}
Las respuestas de las tareas pueden ser enviadas al siguiente correo giovannilopez9808@gmail.com con el asunto entre corchetes de la siguiente manera [Tarea \#] y el nombre del archivo con sus apellidos y nombre, ya sea el nombre completo o solo una parte, por ejemplo \textit{LopezPadilla.pdf}, \textit{LopezGiovanni.pdf}, \textit{LopezPadillaGiovanniGamaliel.pdf}. Las respuestas envienlas en formato pdf, estas pueden ser totalmente digitales (que escriban todo en la computadora) o que sean fotografías, pero por favor, envielo en solo un documento.
\end{document}